\documentclass[11pt,a4paper]{article}
\usepackage[utf8]{inputenc}
\usepackage[T1]{fontenc}
\usepackage[spanish,es-minimal,es-noshorthands]{babel}
\AtBeginDocument{\renewcommand{\sin}{\operatorname{sen}}\renewcommand{\arcsin}{\operatorname{arcsen}}}
\usepackage{lmodern}
\usepackage[sfdefault,lining]{FiraSans}
\usepackage[a4paper,top=2.3cm,bottom=2.2cm,left=1.8cm,right=1.8cm,headheight=1.3cm,headsep=0.4cm,footskip=1cm]{geometry}
\usepackage{amsmath,amssymb}
\usepackage{icomma}
\usepackage{graphicx}
\usepackage[table]{xcolor}
\usepackage{tikz}
\usepackage{fontawesome5}
\usepackage{ccicons}
\usepackage{fancyhdr}
\usepackage{lastpage}
\usepackage{tasks}
\usepackage{needspace}
\usepackage[most]{tcolorbox}
\usepackage{array}
\usepackage[hidelinks]{hyperref}
\AtBeginDocument{\hypersetup{pdftitle={Ficha de ejercicios de la unidad 10: Funciones. 2º ESO},pdfsubject={Matemáticas 2º ESO},pdfkeywords={funciones, variable independiente, variable dependiente, tablas de valores, gráficas, ecuación de una función, función de proporcionalidad, pendiente, función lineal, ordenada en el origen, función constante, interpretación de gráficas}}}

\definecolor{navy}{HTML}{1B2748}
\definecolor{accent}{HTML}{E8603A}
\definecolor{teal}{HTML}{1F9E90}
\definecolor{soft}{HTML}{EEF1F8}
\definecolor{muted}{HTML}{5D6685}

\setlength{\parindent}{0pt}
\setlength{\parskip}{3pt}
\settasks{label=\alph*), label-format=\bfseries\color{navy}, label-width=1.3em, item-indent=1.8em, column-sep=0.6em, after-item-skip=5pt, before-skip=3pt, after-skip=4pt}

% ---------- capçalera i peu ----------
\newcommand{\llicencia}{\href{https://creativecommons.org/licenses/by-nc-sa/4.0/deed.es}{\ccbyncsa}\ \ Licencia Creative Commons Atribución-NoComercial-CompartirIgual 4.0 Internacional (CC BY-NC-SA 4.0)}
\pagestyle{fancy}
\fancyhf{}
\fancyhead[L]{\raisebox{-0.35\height}{\includegraphics[height=0.95cm]{logo-gva.png}}}
\fancyhead[C]{\small\color{muted}Ficha de ejercicios · Unidad 10: Funciones · 2\textsuperscript{o} ESO}
\fancyhead[R]{\raisebox{-0.35\height}{\includegraphics[height=0.95cm]{logo-vermellar.png}}}
\fancyfoot[L]{\scriptsize\color{muted}\llicencia}
\fancyfoot[R]{\small\color{muted}\thepage}
\renewcommand{\headrulewidth}{0.4pt}
\renewcommand{\headrule}{\color{muted!40}\hrule height \headrulewidth}
\fancypagestyle{primera}{\fancyhf{}\renewcommand{\headrulewidth}{0pt}
  \fancyfoot[L]{\scriptsize\color{muted}\llicencia}
  \fancyfoot[R]{\small\color{muted}\thepage}}

% ---------- nivells i exercicis ----------
\newcommand{\estrelles}[1]{{\color{accent}\ifcase#1\or\faStar\faStar[regular]\faStar[regular]\or\faStar\faStar\faStar[regular]\or\faStar\faStar\faStar\fi}}
\newcounter{exer}
\newcommand{\ex}[1]{\par\needspace{4\baselineskip}\medskip\stepcounter{exer}\noindent{\color{navy}\bfseries\theexer.}\ \ {\small\estrelles{#1}}\ \ }

\newcommand{\seccio}[2]{\par\needspace{10\baselineskip}\bigskip
  \begin{tcolorbox}[colback=navy,colframe=navy,coltext=white,arc=2mm,boxrule=0pt,left=3mm,right=3mm,top=1.2mm,bottom=1.2mm]
  {\large\bfseries #1\quad #2}
  \end{tcolorbox}\vspace{-2pt}}

\newcolumntype{C}[1]{>{\centering\arraybackslash}p{#1}}

% quadre de recordatori (només a la fitxa d'exercicis)
\newcommand{\recorda}[1]{\par\smallskip\begin{tcolorbox}[colback=accent!7,colframe=accent!55,arc=2mm,boxrule=0.5pt,left=2.5mm,right=2.5mm,top=1mm,bottom=1mm]\small{\color{accent}\faLightbulb[regular]}\ \textbf{Recuerda:}\ #1\end{tcolorbox}}
\newcommand{\res}[1]{#1}\newcommand{\rest}[1]{#1}

\def\fmneed{14}
% =====================================================================
% Capçalera FentMates de les fitxes: un color per a cada curs.
% Els colors són els de la pàgina inicial (LEVELS a assets/js/materials.js);
% el logo ve de _eines/logo/pdf/ (python3 _eines/logo/genera.py).
%
% Ús, al final del preàmbul de la unitat:
%   % =====================================================================
% Capçalera FentMates de les fitxes: un color per a cada curs.
% Els colors són els de la pàgina inicial (LEVELS a assets/js/materials.js);
% el logo ve de _eines/logo/pdf/ (python3 _eines/logo/genera.py).
%
% Ús, al final del preàmbul de la unitat:
%   % =====================================================================
% Capçalera FentMates de les fitxes: un color per a cada curs.
% Els colors són els de la pàgina inicial (LEVELS a assets/js/materials.js);
% el logo ve de _eines/logo/pdf/ (python3 _eines/logo/genera.py).
%
% Ús, al final del preàmbul de la unitat:
%   \input{../comu/fentmates.tex}
%   \fmcurs{1eso}
% i, al començament del document:
%   \capcalera{Matemàtiques}{1r ESO}{Fitxa d'exercicis · Unitat 2}{Divisibilitat}
% Opcional abans de \input: \def\fmneed{14} (espai mínim abans de cada secció).
% =====================================================================
\makeatletter
\def\fm@logodir{../../logo/pdf}
\providecommand{\fmneed}{10}
\def\fm@cursnom{}

% clau del curs, color i color del text damunt del color
\newcommand{\fm@curs}[3]{\@namedef{fm@color@#1}{#2}\@namedef{fm@text@#1}{#3}}
\fm@curs{1eso}{E6B84A}{navy}
\fm@curs{2eso}{3B82F6}{white}
\fm@curs{3eso}{FF7A50}{white}
\fm@curs{4eso-a}{A78BFA}{white}
\fm@curs{4eso-b}{6D28D9}{white}
\fm@curs{1batx-mi}{9B1C1C}{white}
\fm@curs{1batx-ccss}{EF6B6B}{white}
\fm@curs{2batx-mii}{1E6B47}{white}
\fm@curs{2batx-ccss}{2BB3A3}{white}

\newcommand{\fmcurs}[1]{%
  \@ifundefined{fm@color@#1}{\PackageError{fentmates}{Curs desconegut: #1}{}}{}%
  \edef\fm@tmp{\noexpand\definecolor{curs}{HTML}{\@nameuse{fm@color@#1}}%
               \noexpand\colorlet{curstext}{\@nameuse{fm@text@#1}}}\fm@tmp
  \colorlet{cursfosc}{curs!72!black}%
  \def\fm@logo{\fm@logodir/vertical-#1.pdf}%
  \def\fm@marca{\fm@logodir/marca-#1.pdf}}

% ---------- capçalera de la primera pàgina ----------
% #1 matèria (petit), #2 curs, #3 tipus de document i unitat, #4 títol
\newcommand{\capcalera}[4]{%
  \gdef\fm@cursnom{#2}%
  \par\noindent
  \begin{tcolorbox}[colback=curs,colframe=curs,boxrule=0pt,arc=2.2mm,width=3.5cm,height=2.15cm,
      valign=center,left=3mm,right=2mm,top=1mm,bottom=1mm,nobeforeafter,box align=center]
    \color{curstext}{\footnotesize #1}\par\vspace{1pt}{\fontsize{19}{21}\selectfont\bfseries #2\par}%
  \end{tcolorbox}\hspace{5mm}%
  \begin{minipage}[c]{\dimexpr\linewidth-3.5cm-5mm-2.9cm\relax}
    {\color{cursfosc}\small\bfseries\MakeUppercase{#3}}\par\vspace{2pt}
    {\color{navy}\fontsize{20}{23}\selectfont\bfseries #4\par}
  \end{minipage}\hfill
  \begin{minipage}[c]{2.6cm}\raggedleft\includegraphics[width=2.6cm]{\fm@logo}\end{minipage}\par
  \vspace{2pt}}

% ---------- barres de secció del color del curs ----------
\renewcommand{\seccio}[2]{\par\needspace{\fmneed\baselineskip}\bigskip
  \begin{tcolorbox}[colback=curs,colframe=curs,coltext=curstext,arc=2mm,boxrule=0pt,left=3mm,right=3mm,top=1.2mm,bottom=1.2mm]
  {\large\bfseries #1\quad #2}
  \end{tcolorbox}\vspace{-2pt}}

% ---------- estrelles de dificultat del color del curs (a Batxillerat no se'n posen) ----------
\@ifundefined{estrelles}{}{%
  \renewcommand{\estrelles}[1]{{\color{curs}\ifcase#1\or\faStar\faStar[regular]\faStar[regular]\or\faStar\faStar\faStar[regular]\or\faStar\faStar\faStar\fi}}}

% ---------- capçalera i peu de les pàgines ----------
% primera pàgina: logos institucionals al peu
\fancypagestyle{primera}{\fancyhf{}\renewcommand{\headrulewidth}{0pt}
  \fancyfoot[L]{\raisebox{-0.32\height}{\includegraphics[height=0.78cm]{logo-gva.png}}\hspace{5mm}%
    \raisebox{-0.32\height}{\includegraphics[height=0.78cm]{logo-vermellar.png}}}
  \fancyfoot[R]{\scriptsize\color{muted}\href{https://creativecommons.org/licenses/by-nc-sa/4.0/deed.ca}{\ccbyncsa}\ \ CC BY-NC-SA 4.0 · fentmates.org\qquad\small\thepage}}
% la resta: el símbol a l'esquerra i el curs a la dreta, amb el títol al centre com abans
\fancyhead[L]{\raisebox{-0.3\height}{\includegraphics[height=0.62cm]{\fm@marca}}}
\fancyhead[R]{\small\bfseries\color{cursfosc}\fm@cursnom}
\renewcommand{\headrule}{\color{curs!45}\hrule height \headrulewidth}
\makeatother

%   \fmcurs{1eso}
% i, al començament del document:
%   \capcalera{Matemàtiques}{1r ESO}{Fitxa d'exercicis · Unitat 2}{Divisibilitat}
% Opcional abans de \input: \def\fmneed{14} (espai mínim abans de cada secció).
% =====================================================================
\makeatletter
\def\fm@logodir{../../logo/pdf}
\providecommand{\fmneed}{10}
\def\fm@cursnom{}

% clau del curs, color i color del text damunt del color
\newcommand{\fm@curs}[3]{\@namedef{fm@color@#1}{#2}\@namedef{fm@text@#1}{#3}}
\fm@curs{1eso}{E6B84A}{navy}
\fm@curs{2eso}{3B82F6}{white}
\fm@curs{3eso}{FF7A50}{white}
\fm@curs{4eso-a}{A78BFA}{white}
\fm@curs{4eso-b}{6D28D9}{white}
\fm@curs{1batx-mi}{9B1C1C}{white}
\fm@curs{1batx-ccss}{EF6B6B}{white}
\fm@curs{2batx-mii}{1E6B47}{white}
\fm@curs{2batx-ccss}{2BB3A3}{white}

\newcommand{\fmcurs}[1]{%
  \@ifundefined{fm@color@#1}{\PackageError{fentmates}{Curs desconegut: #1}{}}{}%
  \edef\fm@tmp{\noexpand\definecolor{curs}{HTML}{\@nameuse{fm@color@#1}}%
               \noexpand\colorlet{curstext}{\@nameuse{fm@text@#1}}}\fm@tmp
  \colorlet{cursfosc}{curs!72!black}%
  \def\fm@logo{\fm@logodir/vertical-#1.pdf}%
  \def\fm@marca{\fm@logodir/marca-#1.pdf}}

% ---------- capçalera de la primera pàgina ----------
% #1 matèria (petit), #2 curs, #3 tipus de document i unitat, #4 títol
\newcommand{\capcalera}[4]{%
  \gdef\fm@cursnom{#2}%
  \par\noindent
  \begin{tcolorbox}[colback=curs,colframe=curs,boxrule=0pt,arc=2.2mm,width=3.5cm,height=2.15cm,
      valign=center,left=3mm,right=2mm,top=1mm,bottom=1mm,nobeforeafter,box align=center]
    \color{curstext}{\footnotesize #1}\par\vspace{1pt}{\fontsize{19}{21}\selectfont\bfseries #2\par}%
  \end{tcolorbox}\hspace{5mm}%
  \begin{minipage}[c]{\dimexpr\linewidth-3.5cm-5mm-2.9cm\relax}
    {\color{cursfosc}\small\bfseries\MakeUppercase{#3}}\par\vspace{2pt}
    {\color{navy}\fontsize{20}{23}\selectfont\bfseries #4\par}
  \end{minipage}\hfill
  \begin{minipage}[c]{2.6cm}\raggedleft\includegraphics[width=2.6cm]{\fm@logo}\end{minipage}\par
  \vspace{2pt}}

% ---------- barres de secció del color del curs ----------
\renewcommand{\seccio}[2]{\par\needspace{\fmneed\baselineskip}\bigskip
  \begin{tcolorbox}[colback=curs,colframe=curs,coltext=curstext,arc=2mm,boxrule=0pt,left=3mm,right=3mm,top=1.2mm,bottom=1.2mm]
  {\large\bfseries #1\quad #2}
  \end{tcolorbox}\vspace{-2pt}}

% ---------- estrelles de dificultat del color del curs (a Batxillerat no se'n posen) ----------
\@ifundefined{estrelles}{}{%
  \renewcommand{\estrelles}[1]{{\color{curs}\ifcase#1\or\faStar\faStar[regular]\faStar[regular]\or\faStar\faStar\faStar[regular]\or\faStar\faStar\faStar\fi}}}

% ---------- capçalera i peu de les pàgines ----------
% primera pàgina: logos institucionals al peu
\fancypagestyle{primera}{\fancyhf{}\renewcommand{\headrulewidth}{0pt}
  \fancyfoot[L]{\raisebox{-0.32\height}{\includegraphics[height=0.78cm]{logo-gva.png}}\hspace{5mm}%
    \raisebox{-0.32\height}{\includegraphics[height=0.78cm]{logo-vermellar.png}}}
  \fancyfoot[R]{\scriptsize\color{muted}\href{https://creativecommons.org/licenses/by-nc-sa/4.0/deed.ca}{\ccbyncsa}\ \ CC BY-NC-SA 4.0 · fentmates.org\qquad\small\thepage}}
% la resta: el símbol a l'esquerra i el curs a la dreta, amb el títol al centre com abans
\fancyhead[L]{\raisebox{-0.3\height}{\includegraphics[height=0.62cm]{\fm@marca}}}
\fancyhead[R]{\small\bfseries\color{cursfosc}\fm@cursnom}
\renewcommand{\headrule}{\color{curs!45}\hrule height \headrulewidth}
\makeatother

%   \fmcurs{1eso}
% i, al començament del document:
%   \capcalera{Matemàtiques}{1r ESO}{Fitxa d'exercicis · Unitat 2}{Divisibilitat}
% Opcional abans de \input: \def\fmneed{14} (espai mínim abans de cada secció).
% =====================================================================
\makeatletter
\def\fm@logodir{../../logo/pdf}
\providecommand{\fmneed}{10}
\def\fm@cursnom{}

% clau del curs, color i color del text damunt del color
\newcommand{\fm@curs}[3]{\@namedef{fm@color@#1}{#2}\@namedef{fm@text@#1}{#3}}
\fm@curs{1eso}{E6B84A}{navy}
\fm@curs{2eso}{3B82F6}{white}
\fm@curs{3eso}{FF7A50}{white}
\fm@curs{4eso-a}{A78BFA}{white}
\fm@curs{4eso-b}{6D28D9}{white}
\fm@curs{1batx-mi}{9B1C1C}{white}
\fm@curs{1batx-ccss}{EF6B6B}{white}
\fm@curs{2batx-mii}{1E6B47}{white}
\fm@curs{2batx-ccss}{2BB3A3}{white}

\newcommand{\fmcurs}[1]{%
  \@ifundefined{fm@color@#1}{\PackageError{fentmates}{Curs desconegut: #1}{}}{}%
  \edef\fm@tmp{\noexpand\definecolor{curs}{HTML}{\@nameuse{fm@color@#1}}%
               \noexpand\colorlet{curstext}{\@nameuse{fm@text@#1}}}\fm@tmp
  \colorlet{cursfosc}{curs!72!black}%
  \def\fm@logo{\fm@logodir/vertical-#1.pdf}%
  \def\fm@marca{\fm@logodir/marca-#1.pdf}}

% ---------- capçalera de la primera pàgina ----------
% #1 matèria (petit), #2 curs, #3 tipus de document i unitat, #4 títol
\newcommand{\capcalera}[4]{%
  \gdef\fm@cursnom{#2}%
  \par\noindent
  \begin{tcolorbox}[colback=curs,colframe=curs,boxrule=0pt,arc=2.2mm,width=3.5cm,height=2.15cm,
      valign=center,left=3mm,right=2mm,top=1mm,bottom=1mm,nobeforeafter,box align=center]
    \color{curstext}{\footnotesize #1}\par\vspace{1pt}{\fontsize{19}{21}\selectfont\bfseries #2\par}%
  \end{tcolorbox}\hspace{5mm}%
  \begin{minipage}[c]{\dimexpr\linewidth-3.5cm-5mm-2.9cm\relax}
    {\color{cursfosc}\small\bfseries\MakeUppercase{#3}}\par\vspace{2pt}
    {\color{navy}\fontsize{20}{23}\selectfont\bfseries #4\par}
  \end{minipage}\hfill
  \begin{minipage}[c]{2.6cm}\raggedleft\includegraphics[width=2.6cm]{\fm@logo}\end{minipage}\par
  \vspace{2pt}}

% ---------- barres de secció del color del curs ----------
\renewcommand{\seccio}[2]{\par\needspace{\fmneed\baselineskip}\bigskip
  \begin{tcolorbox}[colback=curs,colframe=curs,coltext=curstext,arc=2mm,boxrule=0pt,left=3mm,right=3mm,top=1.2mm,bottom=1.2mm]
  {\large\bfseries #1\quad #2}
  \end{tcolorbox}\vspace{-2pt}}

% ---------- estrelles de dificultat del color del curs (a Batxillerat no se'n posen) ----------
\@ifundefined{estrelles}{}{%
  \renewcommand{\estrelles}[1]{{\color{curs}\ifcase#1\or\faStar\faStar[regular]\faStar[regular]\or\faStar\faStar\faStar[regular]\or\faStar\faStar\faStar\fi}}}

% ---------- capçalera i peu de les pàgines ----------
% primera pàgina: logos institucionals al peu
\fancypagestyle{primera}{\fancyhf{}\renewcommand{\headrulewidth}{0pt}
  \fancyfoot[L]{\raisebox{-0.32\height}{\includegraphics[height=0.78cm]{logo-gva.png}}\hspace{5mm}%
    \raisebox{-0.32\height}{\includegraphics[height=0.78cm]{logo-vermellar.png}}}
  \fancyfoot[R]{\scriptsize\color{muted}\href{https://creativecommons.org/licenses/by-nc-sa/4.0/deed.ca}{\ccbyncsa}\ \ CC BY-NC-SA 4.0 · fentmates.org\qquad\small\thepage}}
% la resta: el símbol a l'esquerra i el curs a la dreta, amb el títol al centre com abans
\fancyhead[L]{\raisebox{-0.3\height}{\includegraphics[height=0.62cm]{\fm@marca}}}
\fancyhead[R]{\small\bfseries\color{cursfosc}\fm@cursnom}
\renewcommand{\headrule}{\color{curs!45}\hrule height \headrulewidth}
\makeatother

\fmcurs{2eso}

\begin{document}
\thispagestyle{primera}

\capcalera{Matemáticas}{2\textsuperscript{o} ESO}{Ficha de ejercicios · Unidad 10}{Funciones}
\vspace{10pt}
Nombre y apellidos: \rule{6.6cm}{0.4pt}\hfill Grupo: \rule{1.6cm}{0.4pt}\hfill Fecha: \rule{2.2cm}{0.4pt}
\vspace{8pt}

\begin{tcolorbox}[colback=soft,colframe=soft,arc=2mm,left=3mm,right=3mm,top=1.5mm,bottom=1.5mm]
\small Ejercicios graduados por niveles de dificultad:\quad \estrelles{1}\ básico \qquad \estrelles{2}\ medio \qquad \estrelles{3}\ avanzado.\\[2pt] Resuélvelos en el cuaderno mostrando todos los pasos. Para dibujar gráficas usa papel cuadriculado y escribe el nombre y la unidad de cada eje. Las secciones siguen el orden de los apuntes de la unidad (10.1--10.8).
\end{tcolorbox}

% =====================================================
\seccio{1}{Concepto de función}

\ex{1} En cada situación, di cuál es la variable independiente ($x$) y cuál la dependiente ($y$):\begin{tasks}(2)
  \task El precio de una pizza y sus ingredientes.
  \task El área de un círculo y su radio.
  \task Los minutos hablados por teléfono y el precio.
  \task La temperatura de un horno y el tiempo que lleva encendido.
\end{tasks}\recorda{La variable \textbf{independiente} $x$ es la que elegimos o la que pasa sola; la \textbf{dependiente} $y$ es la que depende de la otra. Una relación es una \textbf{función} si a cada valor de $x$ le corresponde \textbf{un único} valor de $y$.}

\ex{1} Di si estas relaciones son funciones y explica por qué:\begin{tasks}(1)
  \task A cada alumno de la clase, su día de nacimiento.
  \task A cada número natural, sus múltiplos.
  \task A cada cantidad de litros de gasolina, el precio que pagamos.
  \task A cada edad, las personas de la clase que la tienen.
\end{tasks}

\ex{1} ¿Cuáles de estas tablas son de una función? Razónalo.\par\smallskip\begin{center}$\begin{array}{|c|cccc|}\hline x&1&2&3&4\\\hline y&3&6&9&12\\\hline\end{array}$\hspace{0.7cm}$\begin{array}{|c|cccc|}\hline x&1&1&2&3\\\hline y&2&5&4&7\\\hline\end{array}$\hspace{0.7cm}$\begin{array}{|c|cccc|}\hline x&-2&-1&1&2\\\hline y&4&1&1&4\\\hline\end{array}$\end{center}

\ex{2} ¿Cuáles de estas gráficas son de una función? Haz la prueba de la recta vertical.\par\smallskip\begin{center}\begin{tabular}{c}\begin{tikzpicture}[x=0.3cm,y=0.3cm,baseline=(current bounding box.center),line cap=round,line join=round]
\draw[muted!30,line width=0.3pt] (-4,-3)--(-4,3);
\draw[muted!30,line width=0.3pt] (-3,-3)--(-3,3);
\draw[muted!30,line width=0.3pt] (-2,-3)--(-2,3);
\draw[muted!30,line width=0.3pt] (-1,-3)--(-1,3);
\draw[muted!30,line width=0.3pt] (0,-3)--(0,3);
\draw[muted!30,line width=0.3pt] (1,-3)--(1,3);
\draw[muted!30,line width=0.3pt] (2,-3)--(2,3);
\draw[muted!30,line width=0.3pt] (3,-3)--(3,3);
\draw[muted!30,line width=0.3pt] (4,-3)--(4,3);
\draw[muted!30,line width=0.3pt] (-4,-3)--(4,-3);
\draw[muted!30,line width=0.3pt] (-4,-2)--(4,-2);
\draw[muted!30,line width=0.3pt] (-4,-1)--(4,-1);
\draw[muted!30,line width=0.3pt] (-4,0)--(4,0);
\draw[muted!30,line width=0.3pt] (-4,1)--(4,1);
\draw[muted!30,line width=0.3pt] (-4,2)--(4,2);
\draw[muted!30,line width=0.3pt] (-4,3)--(4,3);
\draw[-latex,muted,line width=0.6pt] (-4,0)--(4.3,0) node[right,font=\scriptsize,text=navy]{x};
\draw[-latex,muted,line width=0.6pt] (0,-3)--(0,3.3) node[above,font=\scriptsize,text=navy]{y};
\node[font=\tiny,text=muted,below,inner sep=1.5pt] at (-2,0) {$-2$};
\node[font=\tiny,text=muted,below,inner sep=1.5pt] at (2,0) {$2$};
\node[font=\tiny,text=muted,left,inner sep=1.5pt] at (0,-2) {$-2$};
\node[font=\tiny,text=muted,left,inner sep=1.5pt] at (0,2) {$2$};
\node[font=\tiny,text=muted,below left,inner sep=1.5pt] at (0,0) {$0$};
\draw[accent,line width=1.1pt] (-3.500,-2.000)--(-1.500,1.500)--(0.500,0.000)--(3.500,2.500);
\end{tikzpicture}\\[1pt]\small A\end{tabular}\hspace{0.7cm}\begin{tabular}{c}\begin{tikzpicture}[x=0.3cm,y=0.3cm,baseline=(current bounding box.center),line cap=round,line join=round]
\draw[muted!30,line width=0.3pt] (-4,-3)--(-4,3);
\draw[muted!30,line width=0.3pt] (-3,-3)--(-3,3);
\draw[muted!30,line width=0.3pt] (-2,-3)--(-2,3);
\draw[muted!30,line width=0.3pt] (-1,-3)--(-1,3);
\draw[muted!30,line width=0.3pt] (0,-3)--(0,3);
\draw[muted!30,line width=0.3pt] (1,-3)--(1,3);
\draw[muted!30,line width=0.3pt] (2,-3)--(2,3);
\draw[muted!30,line width=0.3pt] (3,-3)--(3,3);
\draw[muted!30,line width=0.3pt] (4,-3)--(4,3);
\draw[muted!30,line width=0.3pt] (-4,-3)--(4,-3);
\draw[muted!30,line width=0.3pt] (-4,-2)--(4,-2);
\draw[muted!30,line width=0.3pt] (-4,-1)--(4,-1);
\draw[muted!30,line width=0.3pt] (-4,0)--(4,0);
\draw[muted!30,line width=0.3pt] (-4,1)--(4,1);
\draw[muted!30,line width=0.3pt] (-4,2)--(4,2);
\draw[muted!30,line width=0.3pt] (-4,3)--(4,3);
\draw[-latex,muted,line width=0.6pt] (-4,0)--(4.3,0) node[right,font=\scriptsize,text=navy]{x};
\draw[-latex,muted,line width=0.6pt] (0,-3)--(0,3.3) node[above,font=\scriptsize,text=navy]{y};
\node[font=\tiny,text=muted,below,inner sep=1.5pt] at (-2,0) {$-2$};
\node[font=\tiny,text=muted,below,inner sep=1.5pt] at (2,0) {$2$};
\node[font=\tiny,text=muted,left,inner sep=1.5pt] at (0,-2) {$-2$};
\node[font=\tiny,text=muted,left,inner sep=1.5pt] at (0,2) {$2$};
\node[font=\tiny,text=muted,below left,inner sep=1.5pt] at (0,0) {$0$};
\draw[accent,line width=1.1pt] (2.200,0.000)--(2.167,0.382)--(2.067,0.752)--(1.905,1.100)--(1.685,1.414)--(1.414,1.685)--(1.100,1.905)--(0.752,2.067)--(0.382,2.167)--(0.000,2.200)--(-0.382,2.167)--(-0.752,2.067)--(-1.100,1.905)--(-1.414,1.685)--(-1.685,1.414)--(-1.905,1.100)--(-2.067,0.752)--(-2.167,0.382)--(-2.200,0.000)--(-2.167,-0.382)--(-2.067,-0.752)--(-1.905,-1.100)--(-1.685,-1.414)--(-1.414,-1.685)--(-1.100,-1.905)--(-0.752,-2.067)--(-0.382,-2.167)--(-0.000,-2.200)--(0.382,-2.167)--(0.752,-2.067)--(1.100,-1.905)--(1.414,-1.685)--(1.685,-1.414)--(1.905,-1.100)--(2.067,-0.752)--(2.167,-0.382)--(2.200,-0.000);
\end{tikzpicture}\\[1pt]\small B\end{tabular}\hspace{0.7cm}\begin{tabular}{c}\begin{tikzpicture}[x=0.3cm,y=0.3cm,baseline=(current bounding box.center),line cap=round,line join=round]
\draw[muted!30,line width=0.3pt] (-4,-3)--(-4,3);
\draw[muted!30,line width=0.3pt] (-3,-3)--(-3,3);
\draw[muted!30,line width=0.3pt] (-2,-3)--(-2,3);
\draw[muted!30,line width=0.3pt] (-1,-3)--(-1,3);
\draw[muted!30,line width=0.3pt] (0,-3)--(0,3);
\draw[muted!30,line width=0.3pt] (1,-3)--(1,3);
\draw[muted!30,line width=0.3pt] (2,-3)--(2,3);
\draw[muted!30,line width=0.3pt] (3,-3)--(3,3);
\draw[muted!30,line width=0.3pt] (4,-3)--(4,3);
\draw[muted!30,line width=0.3pt] (-4,-3)--(4,-3);
\draw[muted!30,line width=0.3pt] (-4,-2)--(4,-2);
\draw[muted!30,line width=0.3pt] (-4,-1)--(4,-1);
\draw[muted!30,line width=0.3pt] (-4,0)--(4,0);
\draw[muted!30,line width=0.3pt] (-4,1)--(4,1);
\draw[muted!30,line width=0.3pt] (-4,2)--(4,2);
\draw[muted!30,line width=0.3pt] (-4,3)--(4,3);
\draw[-latex,muted,line width=0.6pt] (-4,0)--(4.3,0) node[right,font=\scriptsize,text=navy]{x};
\draw[-latex,muted,line width=0.6pt] (0,-3)--(0,3.3) node[above,font=\scriptsize,text=navy]{y};
\node[font=\tiny,text=muted,below,inner sep=1.5pt] at (-2,0) {$-2$};
\node[font=\tiny,text=muted,below,inner sep=1.5pt] at (2,0) {$2$};
\node[font=\tiny,text=muted,left,inner sep=1.5pt] at (0,-2) {$-2$};
\node[font=\tiny,text=muted,left,inner sep=1.5pt] at (0,2) {$2$};
\node[font=\tiny,text=muted,below left,inner sep=1.5pt] at (0,0) {$0$};
\draw[accent,line width=1.1pt] (3.750,-2.500)--(3.260,-2.400)--(2.790,-2.300)--(2.340,-2.200)--(1.910,-2.100)--(1.500,-2.000)--(1.110,-1.900)--(0.740,-1.800)--(0.390,-1.700)--(0.060,-1.600)--(-0.250,-1.500)--(-0.540,-1.400)--(-0.810,-1.300)--(-1.060,-1.200)--(-1.290,-1.100)--(-1.500,-1.000)--(-1.690,-0.900)--(-1.860,-0.800)--(-2.010,-0.700)--(-2.140,-0.600)--(-2.250,-0.500)--(-2.340,-0.400)--(-2.410,-0.300)--(-2.460,-0.200)--(-2.490,-0.100)--(-2.500,0.000)--(-2.490,0.100)--(-2.460,0.200)--(-2.410,0.300)--(-2.340,0.400)--(-2.250,0.500)--(-2.140,0.600)--(-2.010,0.700)--(-1.860,0.800)--(-1.690,0.900)--(-1.500,1.000)--(-1.290,1.100)--(-1.060,1.200)--(-0.810,1.300)--(-0.540,1.400)--(-0.250,1.500)--(0.060,1.600)--(0.390,1.700)--(0.740,1.800)--(1.110,1.900)--(1.500,2.000)--(1.910,2.100)--(2.340,2.200)--(2.790,2.300)--(3.260,2.400)--(3.750,2.500);
\end{tikzpicture}\\[1pt]\small C\end{tabular}\hspace{0.7cm}\begin{tabular}{c}\begin{tikzpicture}[x=0.3cm,y=0.3cm,baseline=(current bounding box.center),line cap=round,line join=round]
\draw[muted!30,line width=0.3pt] (-4,-3)--(-4,3);
\draw[muted!30,line width=0.3pt] (-3,-3)--(-3,3);
\draw[muted!30,line width=0.3pt] (-2,-3)--(-2,3);
\draw[muted!30,line width=0.3pt] (-1,-3)--(-1,3);
\draw[muted!30,line width=0.3pt] (0,-3)--(0,3);
\draw[muted!30,line width=0.3pt] (1,-3)--(1,3);
\draw[muted!30,line width=0.3pt] (2,-3)--(2,3);
\draw[muted!30,line width=0.3pt] (3,-3)--(3,3);
\draw[muted!30,line width=0.3pt] (4,-3)--(4,3);
\draw[muted!30,line width=0.3pt] (-4,-3)--(4,-3);
\draw[muted!30,line width=0.3pt] (-4,-2)--(4,-2);
\draw[muted!30,line width=0.3pt] (-4,-1)--(4,-1);
\draw[muted!30,line width=0.3pt] (-4,0)--(4,0);
\draw[muted!30,line width=0.3pt] (-4,1)--(4,1);
\draw[muted!30,line width=0.3pt] (-4,2)--(4,2);
\draw[muted!30,line width=0.3pt] (-4,3)--(4,3);
\draw[-latex,muted,line width=0.6pt] (-4,0)--(4.3,0) node[right,font=\scriptsize,text=navy]{x};
\draw[-latex,muted,line width=0.6pt] (0,-3)--(0,3.3) node[above,font=\scriptsize,text=navy]{y};
\node[font=\tiny,text=muted,below,inner sep=1.5pt] at (-2,0) {$-2$};
\node[font=\tiny,text=muted,below,inner sep=1.5pt] at (2,0) {$2$};
\node[font=\tiny,text=muted,left,inner sep=1.5pt] at (0,-2) {$-2$};
\node[font=\tiny,text=muted,left,inner sep=1.5pt] at (0,2) {$2$};
\node[font=\tiny,text=muted,below left,inner sep=1.5pt] at (0,0) {$0$};
\draw[accent,line width=1.1pt] (-4.000,-1.500)--(4.000,-1.500);
\end{tikzpicture}\\[1pt]\small D\end{tabular}\end{center}\recorda{Si alguna recta vertical corta la gráfica en dos puntos o más, \textbf{no} es una función.}

\ex{2} Una impresora imprime 20 páginas cada minuto. Da esta función con una tabla (de 0 a 4 minutos), con una fórmula y con una gráfica. ¿Unimos los puntos?

\ex{3} Los rectángulos de área 24\,cm$^2$ pueden tener bases distintas. a) Completa la tabla de la altura $h$ según la base $b$. b) ¿Es una función? c) ¿Qué fórmula tiene? d) Cuando la base se duplica, ¿qué pasa con la altura?\par\smallskip\begin{center}$\begin{array}{|c|C{0.75cm}C{0.75cm}C{0.75cm}C{0.75cm}C{0.75cm}C{0.75cm}C{0.75cm}C{0.75cm}|}\hline \text{b}&1&2&3&4&6&8&12&24\\\hline \text{h}&\phantom{00}&\phantom{00}&\phantom{00}&\phantom{00}&\phantom{00}&\phantom{00}&\phantom{00}&\phantom{00}\\\hline\end{array}$\end{center}


% =====================================================
\seccio{2}{Funciones dadas por tablas de valores}

\ex{1} Las mandarinas valen 1,20\,€ el kilo. Completa la tabla del precio ($y$) según el peso ($x$) y di si en la gráfica unimos los puntos.\par\smallskip\begin{center}$\begin{array}{|c|C{0.75cm}C{0.75cm}C{0.75cm}C{0.75cm}C{0.75cm}|}\hline \text{kg}&0,5&1&2&3&4\\\hline \text{€}&\phantom{00}&\phantom{00}&\phantom{00}&\phantom{00}&\phantom{00}\\\hline\end{array}$\end{center}\recorda{Cada columna de la tabla es un punto $(x,\ y)$ de la gráfica. Unimos los puntos si también tienen sentido los valores intermedios; si solo hay valores sueltos (personas, entradas…), no los unimos.}

\ex{1} Di si en la gráfica de estas funciones se tienen que unir los puntos:\begin{tasks}(1)
  \task Número de hermanos de un alumno → número de regalos que hace en Navidad.
  \task Tiempo que lleva corriendo un atleta → distancia recorrida.
  \task Número de libros comprados → precio total.
  \task Litros de agua de un depósito → altura del agua.
\end{tasks}

\ex{1} La tabla da la temperatura de un día cada 3 horas. Represéntala y responde: a) ¿A qué hora hizo más calor? ¿Y más frío? b) ¿Entre qué horas subió la temperatura? c) ¿Qué temperatura aproximada hacía a las 10 h?\par\smallskip\begin{center}$\begin{array}{|c|ccccccccc|}\hline \text{h}&0&3&6&9&12&15&18&21&24\\\hline \text{°C}&8&6&7&13&18&20&16&11&9\\\hline\end{array}$\end{center}

\ex{2} Halla la fórmula de la función que da cada tabla:\par\smallskip\begin{center}$\begin{array}{|c|ccccc|}\hline x&0&1&2&3&4\\\hline y&1&4&7&10&13\\\hline\end{array}$\hspace{0.7cm}$\begin{array}{|c|ccccc|}\hline x&0&1&2&3&4\\\hline y&20&17&14&11&8\\\hline\end{array}$\end{center}\recorda{Si $x$ va de 1 en 1 e $y$ aumenta (o disminuye) siempre igual, la fórmula es $y=mx+n$: $m$ es lo que cambia $y$ cada vez y $n$ es el valor de $y$ para $x=0$.}

\ex{2} Completa la tabla de la función $y=2x-5$:\par\smallskip\begin{center}$\begin{array}{|c|cccccc|}\hline x&-1&0&2&\phantom{00}&\phantom{00}&10\\\hline y&\phantom{00}&\phantom{00}&\phantom{00}&3&0&\phantom{00}\\\hline\end{array}$\end{center}

\ex{2} Una piscina se vacía. La tabla da la altura del agua (cm) según las horas. a) ¿Cuánta agua baja cada hora? b) ¿Qué altura había a las 5 h? c) ¿Cuándo quedarán 20\,cm? d) Escribe la fórmula.\par\smallskip\begin{center}$\begin{array}{|c|ccccc|}\hline \text{h}&0&2&4&6&8\\\hline \text{cm}&90&70&50&30&10\\\hline\end{array}$\end{center}

\ex{3} Haz la tabla del área $A$ de un cuadrado según el lado $l$ para $l=1,\ 2,\ 3,\ 4,\ 5$ y represéntala. ¿Es una recta? ¿Qué área tiene un cuadrado de 2,5\,cm de lado?

\ex{3} Con palillos hacemos una fila de triángulos:\par\smallskip\begin{center}\begin{tikzpicture}[x=0.5cm,y=0.5cm,line width=1pt,accent]\draw (0,0)--(0.5,0.87)--(1,0);\draw (0.5,0.87)--(1.5,0.87);\draw (1,0)--(1.5,0.87)--(2,0);\draw (1.5,0.87)--(2.5,0.87);\draw (2,0)--(2.5,0.87)--(3,0);\draw (2.5,0.87)--(3.5,0.87);\draw (3,0)--(3.5,0.87)--(4,0);\end{tikzpicture}\end{center}a) Completa la tabla del número de palillos $y$ según el número de triángulos $x$. b) Halla su fórmula. c) ¿Cuántos palillos hacen falta para 20 triángulos? d) ¿Cuántos triángulos hacemos con 61 palillos?\par\smallskip\begin{center}$\begin{array}{|c|C{0.75cm}C{0.75cm}C{0.75cm}C{0.75cm}|}\hline \text{triángulos}&1&2&3&4\\\hline \text{palillos}&\phantom{00}&\phantom{00}&\phantom{00}&\phantom{00}\\\hline\end{array}$\end{center}


% =====================================================
\seccio{3}{Funciones dadas por su ecuación}

\ex{1} Calcula el valor de $y=4x-3$ para:\begin{tasks}(4)
  \task $x=-2$
  \task $x=0$
  \task $x=\frac12$
  \task $x=3$
\end{tasks}\recorda{Sustituimos la $x$ por su valor (entre paréntesis si es negativo) y hacemos las operaciones. El punto $(x,\ y)$ obtenido es de la gráfica.}

\ex{1} Haz una tabla de valores de $y=-x+2$ para $x$ de $-2$ a $4$ y represéntala.

\ex{1} Haz la tabla de valores de $y=x^2-2$ para $x=-3,\ -2,\ \ldots,\ 3$ y dibuja su gráfica.

\ex{2} ¿Cuáles de estos puntos son de la gráfica de $y=3x-1$?\begin{tasks}(4)
  \task $A(2,\ 5)$
  \task $B(-1,\ -4)$
  \task $C(0,\ 1)$
  \task $D\left(\frac13,\ 0\right)$
\end{tasks}

\ex{2} Halla el valor de $x$ en cada caso:\begin{tasks}(1)
  \task ¿Para qué valor de $x$ la función $y=5x+2$ vale 17?
  \task ¿Para qué valor de $x$ la función $y=-2x+3$ vale 9?
  \task ¿Para qué valores de $x$ la función $y=x^2$ vale 16?
\end{tasks}

\ex{2} Haz la tabla de $y=\frac{12}{x}$ para $x=1,\ 2,\ 3,\ 4,\ 6,\ 12$ y dibuja su gráfica. ¿Cuánto vale $y$ para $x=0$?

\ex{3} Un rectángulo tiene 16\,m de perímetro y base $x$. a) Comprueba que la altura es $8-x$ y escribe la ecuación del área $A$. b) Haz la tabla para $x=1,\ \ldots,\ 7$ y la gráfica. c) ¿Qué dimensiones dan el área máxima?

\ex{3} La distancia que necesita un coche para frenar (en m) se puede calcular con $d=\frac{v^2}{100}$, donde $v$ es la velocidad en km/h. a) Haz la tabla para 30, 50, 80, 100 y 120\,km/h. b) Si la velocidad se duplica de 50 a 100\,km/h, ¿la distancia también se duplica? c) ¿Qué distancia necesita a 90\,km/h?


% =====================================================
\seccio{4}{Función de proporcionalidad: $y=mx$}

\ex{1} ¿Qué tablas corresponden a una función de proporcionalidad? Escribe su ecuación.\par\smallskip\begin{center}$\begin{array}{|c|ccc|}\hline x&1&2&4\\\hline y&5&10&20\\\hline\end{array}$\hspace{0.7cm}$\begin{array}{|c|ccc|}\hline x&2&3&5\\\hline y&6&8&12\\\hline\end{array}$\hspace{0.7cm}$\begin{array}{|c|ccc|}\hline x&2&4&6\\\hline y&-1&-2&-3\\\hline\end{array}$\end{center}\recorda{En una función de proporcionalidad $y=mx$, el cociente $\frac yx$ es siempre $m$. La gráfica es una recta que pasa por el origen.}

\ex{1} Representa en los mismos ejes $y=2x$, $y=-3x$ e $y=\frac x3$. ¿Cuáles son crecientes? ¿Cuál es la más inclinada?

\ex{1} Escribe la ecuación de la recta que pasa por el origen y por el punto:\begin{tasks}(3)
  \task $(3,\ 12)$
  \task $(-2,\ 5)$
  \task $(4,\ -2)$
\end{tasks}

\ex{2} Una tela cuesta 4,50\,€ el metro. Escribe la función del precio $y$ según los metros $x$, haz una tabla y la gráfica. ¿Cuántos metros compramos con 27\,€?

\ex{2} Un coche gasta 6 litros de gasolina cada 100\,km. a) Escribe la función de los litros $y$ que gasta según los km $x$. b) ¿Cuánto gasta en 350\,km? c) ¿Cuántos km puede hacer con 45 litros?

\ex{2} Sin dibujarlas, di si estas rectas son crecientes o decrecientes y ordénalas de menos a más inclinada: $y=-5x$, \ $y=\frac12x$, \ $y=2x$, \ $y=-x$.

\ex{3} Un día, 1 euro equivale a 0,85 libras. a) Escribe la función que pasa de euros ($x$) a libras ($y$). b) ¿Cuántas libras son 60\,€? c) ¿Cuántos euros son 34 libras? d) Escribe también la función que pasa de libras a euros. ¿Es de proporcionalidad?


% =====================================================
\seccio{5}{Pendiente de una recta}

\ex{1} Calcula la pendiente de la recta que pasa por cada par de puntos y di si es creciente, decreciente o constante:\begin{tasks}(2)
  \task $(1,\ 2)$ y $(4,\ 11)$
  \task $(-2,\ 5)$ y $(1,\ -1)$
  \task $(0,\ 3)$ y $(6,\ 7)$
  \task $(-3,\ 4)$ y $(5,\ 4)$
\end{tasks}\recorda{$m=\dfrac{y_2-y_1}{x_2-x_1}$: es lo que aumenta $y$ cuando $x$ aumenta 1. Si $m>0$ la recta sube; si $m<0$, baja; si $m=0$, es horizontal.}

\ex{1} Lee la pendiente de cada recta en la gráfica (busca dos puntos que caigan en la cuadrícula):\par\smallskip\begin{center}\begin{tabular}{c}\begin{tikzpicture}[x=0.3cm,y=0.3cm,baseline=(current bounding box.center),line cap=round,line join=round]
\draw[muted!30,line width=0.3pt] (-4,-4)--(-4,4);
\draw[muted!30,line width=0.3pt] (-3,-4)--(-3,4);
\draw[muted!30,line width=0.3pt] (-2,-4)--(-2,4);
\draw[muted!30,line width=0.3pt] (-1,-4)--(-1,4);
\draw[muted!30,line width=0.3pt] (0,-4)--(0,4);
\draw[muted!30,line width=0.3pt] (1,-4)--(1,4);
\draw[muted!30,line width=0.3pt] (2,-4)--(2,4);
\draw[muted!30,line width=0.3pt] (3,-4)--(3,4);
\draw[muted!30,line width=0.3pt] (4,-4)--(4,4);
\draw[muted!30,line width=0.3pt] (-4,-4)--(4,-4);
\draw[muted!30,line width=0.3pt] (-4,-3)--(4,-3);
\draw[muted!30,line width=0.3pt] (-4,-2)--(4,-2);
\draw[muted!30,line width=0.3pt] (-4,-1)--(4,-1);
\draw[muted!30,line width=0.3pt] (-4,0)--(4,0);
\draw[muted!30,line width=0.3pt] (-4,1)--(4,1);
\draw[muted!30,line width=0.3pt] (-4,2)--(4,2);
\draw[muted!30,line width=0.3pt] (-4,3)--(4,3);
\draw[muted!30,line width=0.3pt] (-4,4)--(4,4);
\draw[-latex,muted,line width=0.6pt] (-4,0)--(4.3,0) node[right,font=\scriptsize,text=navy]{x};
\draw[-latex,muted,line width=0.6pt] (0,-4)--(0,4.3) node[above,font=\scriptsize,text=navy]{y};
\node[font=\tiny,text=muted,below,inner sep=1.5pt] at (-3,0) {$-3$};
\node[font=\tiny,text=muted,below,inner sep=1.5pt] at (-2,0) {$-2$};
\node[font=\tiny,text=muted,below,inner sep=1.5pt] at (-1,0) {$-1$};
\node[font=\tiny,text=muted,below,inner sep=1.5pt] at (1,0) {$1$};
\node[font=\tiny,text=muted,below,inner sep=1.5pt] at (2,0) {$2$};
\node[font=\tiny,text=muted,below,inner sep=1.5pt] at (3,0) {$3$};
\node[font=\tiny,text=muted,left,inner sep=1.5pt] at (0,-3) {$-3$};
\node[font=\tiny,text=muted,left,inner sep=1.5pt] at (0,-2) {$-2$};
\node[font=\tiny,text=muted,left,inner sep=1.5pt] at (0,-1) {$-1$};
\node[font=\tiny,text=muted,left,inner sep=1.5pt] at (0,1) {$1$};
\node[font=\tiny,text=muted,left,inner sep=1.5pt] at (0,2) {$2$};
\node[font=\tiny,text=muted,left,inner sep=1.5pt] at (0,3) {$3$};
\node[font=\tiny,text=muted,below left,inner sep=1.5pt] at (0,0) {$0$};
\draw[accent,line width=1.1pt] (-1.500,-4.000)--(2.500,4.000);
\end{tikzpicture}\\[1pt]\small a)\end{tabular}\hspace{0.7cm}\begin{tabular}{c}\begin{tikzpicture}[x=0.3cm,y=0.3cm,baseline=(current bounding box.center),line cap=round,line join=round]
\draw[muted!30,line width=0.3pt] (-4,-4)--(-4,4);
\draw[muted!30,line width=0.3pt] (-3,-4)--(-3,4);
\draw[muted!30,line width=0.3pt] (-2,-4)--(-2,4);
\draw[muted!30,line width=0.3pt] (-1,-4)--(-1,4);
\draw[muted!30,line width=0.3pt] (0,-4)--(0,4);
\draw[muted!30,line width=0.3pt] (1,-4)--(1,4);
\draw[muted!30,line width=0.3pt] (2,-4)--(2,4);
\draw[muted!30,line width=0.3pt] (3,-4)--(3,4);
\draw[muted!30,line width=0.3pt] (4,-4)--(4,4);
\draw[muted!30,line width=0.3pt] (-4,-4)--(4,-4);
\draw[muted!30,line width=0.3pt] (-4,-3)--(4,-3);
\draw[muted!30,line width=0.3pt] (-4,-2)--(4,-2);
\draw[muted!30,line width=0.3pt] (-4,-1)--(4,-1);
\draw[muted!30,line width=0.3pt] (-4,0)--(4,0);
\draw[muted!30,line width=0.3pt] (-4,1)--(4,1);
\draw[muted!30,line width=0.3pt] (-4,2)--(4,2);
\draw[muted!30,line width=0.3pt] (-4,3)--(4,3);
\draw[muted!30,line width=0.3pt] (-4,4)--(4,4);
\draw[-latex,muted,line width=0.6pt] (-4,0)--(4.3,0) node[right,font=\scriptsize,text=navy]{x};
\draw[-latex,muted,line width=0.6pt] (0,-4)--(0,4.3) node[above,font=\scriptsize,text=navy]{y};
\node[font=\tiny,text=muted,below,inner sep=1.5pt] at (-3,0) {$-3$};
\node[font=\tiny,text=muted,below,inner sep=1.5pt] at (-2,0) {$-2$};
\node[font=\tiny,text=muted,below,inner sep=1.5pt] at (-1,0) {$-1$};
\node[font=\tiny,text=muted,below,inner sep=1.5pt] at (1,0) {$1$};
\node[font=\tiny,text=muted,below,inner sep=1.5pt] at (2,0) {$2$};
\node[font=\tiny,text=muted,below,inner sep=1.5pt] at (3,0) {$3$};
\node[font=\tiny,text=muted,left,inner sep=1.5pt] at (0,-3) {$-3$};
\node[font=\tiny,text=muted,left,inner sep=1.5pt] at (0,-2) {$-2$};
\node[font=\tiny,text=muted,left,inner sep=1.5pt] at (0,-1) {$-1$};
\node[font=\tiny,text=muted,left,inner sep=1.5pt] at (0,1) {$1$};
\node[font=\tiny,text=muted,left,inner sep=1.5pt] at (0,2) {$2$};
\node[font=\tiny,text=muted,left,inner sep=1.5pt] at (0,3) {$3$};
\node[font=\tiny,text=muted,below left,inner sep=1.5pt] at (0,0) {$0$};
\draw[accent,line width=1.1pt] (-4.000,4.000)--(4.000,0.000);
\end{tikzpicture}\\[1pt]\small b)\end{tabular}\hspace{0.7cm}\begin{tabular}{c}\begin{tikzpicture}[x=0.3cm,y=0.3cm,baseline=(current bounding box.center),line cap=round,line join=round]
\draw[muted!30,line width=0.3pt] (-4,-4)--(-4,4);
\draw[muted!30,line width=0.3pt] (-3,-4)--(-3,4);
\draw[muted!30,line width=0.3pt] (-2,-4)--(-2,4);
\draw[muted!30,line width=0.3pt] (-1,-4)--(-1,4);
\draw[muted!30,line width=0.3pt] (0,-4)--(0,4);
\draw[muted!30,line width=0.3pt] (1,-4)--(1,4);
\draw[muted!30,line width=0.3pt] (2,-4)--(2,4);
\draw[muted!30,line width=0.3pt] (3,-4)--(3,4);
\draw[muted!30,line width=0.3pt] (4,-4)--(4,4);
\draw[muted!30,line width=0.3pt] (-4,-4)--(4,-4);
\draw[muted!30,line width=0.3pt] (-4,-3)--(4,-3);
\draw[muted!30,line width=0.3pt] (-4,-2)--(4,-2);
\draw[muted!30,line width=0.3pt] (-4,-1)--(4,-1);
\draw[muted!30,line width=0.3pt] (-4,0)--(4,0);
\draw[muted!30,line width=0.3pt] (-4,1)--(4,1);
\draw[muted!30,line width=0.3pt] (-4,2)--(4,2);
\draw[muted!30,line width=0.3pt] (-4,3)--(4,3);
\draw[muted!30,line width=0.3pt] (-4,4)--(4,4);
\draw[-latex,muted,line width=0.6pt] (-4,0)--(4.3,0) node[right,font=\scriptsize,text=navy]{x};
\draw[-latex,muted,line width=0.6pt] (0,-4)--(0,4.3) node[above,font=\scriptsize,text=navy]{y};
\node[font=\tiny,text=muted,below,inner sep=1.5pt] at (-3,0) {$-3$};
\node[font=\tiny,text=muted,below,inner sep=1.5pt] at (-2,0) {$-2$};
\node[font=\tiny,text=muted,below,inner sep=1.5pt] at (-1,0) {$-1$};
\node[font=\tiny,text=muted,below,inner sep=1.5pt] at (1,0) {$1$};
\node[font=\tiny,text=muted,below,inner sep=1.5pt] at (2,0) {$2$};
\node[font=\tiny,text=muted,below,inner sep=1.5pt] at (3,0) {$3$};
\node[font=\tiny,text=muted,left,inner sep=1.5pt] at (0,-3) {$-3$};
\node[font=\tiny,text=muted,left,inner sep=1.5pt] at (0,-2) {$-2$};
\node[font=\tiny,text=muted,left,inner sep=1.5pt] at (0,-1) {$-1$};
\node[font=\tiny,text=muted,left,inner sep=1.5pt] at (0,1) {$1$};
\node[font=\tiny,text=muted,left,inner sep=1.5pt] at (0,2) {$2$};
\node[font=\tiny,text=muted,left,inner sep=1.5pt] at (0,3) {$3$};
\node[font=\tiny,text=muted,below left,inner sep=1.5pt] at (0,0) {$0$};
\draw[accent,line width=1.1pt] (-4.000,-4.000)--(4.000,2.000);
\end{tikzpicture}\\[1pt]\small c)\end{tabular}\hspace{0.7cm}\begin{tabular}{c}\begin{tikzpicture}[x=0.3cm,y=0.3cm,baseline=(current bounding box.center),line cap=round,line join=round]
\draw[muted!30,line width=0.3pt] (-4,-4)--(-4,4);
\draw[muted!30,line width=0.3pt] (-3,-4)--(-3,4);
\draw[muted!30,line width=0.3pt] (-2,-4)--(-2,4);
\draw[muted!30,line width=0.3pt] (-1,-4)--(-1,4);
\draw[muted!30,line width=0.3pt] (0,-4)--(0,4);
\draw[muted!30,line width=0.3pt] (1,-4)--(1,4);
\draw[muted!30,line width=0.3pt] (2,-4)--(2,4);
\draw[muted!30,line width=0.3pt] (3,-4)--(3,4);
\draw[muted!30,line width=0.3pt] (4,-4)--(4,4);
\draw[muted!30,line width=0.3pt] (-4,-4)--(4,-4);
\draw[muted!30,line width=0.3pt] (-4,-3)--(4,-3);
\draw[muted!30,line width=0.3pt] (-4,-2)--(4,-2);
\draw[muted!30,line width=0.3pt] (-4,-1)--(4,-1);
\draw[muted!30,line width=0.3pt] (-4,0)--(4,0);
\draw[muted!30,line width=0.3pt] (-4,1)--(4,1);
\draw[muted!30,line width=0.3pt] (-4,2)--(4,2);
\draw[muted!30,line width=0.3pt] (-4,3)--(4,3);
\draw[muted!30,line width=0.3pt] (-4,4)--(4,4);
\draw[-latex,muted,line width=0.6pt] (-4,0)--(4.3,0) node[right,font=\scriptsize,text=navy]{x};
\draw[-latex,muted,line width=0.6pt] (0,-4)--(0,4.3) node[above,font=\scriptsize,text=navy]{y};
\node[font=\tiny,text=muted,below,inner sep=1.5pt] at (-3,0) {$-3$};
\node[font=\tiny,text=muted,below,inner sep=1.5pt] at (-2,0) {$-2$};
\node[font=\tiny,text=muted,below,inner sep=1.5pt] at (-1,0) {$-1$};
\node[font=\tiny,text=muted,below,inner sep=1.5pt] at (1,0) {$1$};
\node[font=\tiny,text=muted,below,inner sep=1.5pt] at (2,0) {$2$};
\node[font=\tiny,text=muted,below,inner sep=1.5pt] at (3,0) {$3$};
\node[font=\tiny,text=muted,left,inner sep=1.5pt] at (0,-3) {$-3$};
\node[font=\tiny,text=muted,left,inner sep=1.5pt] at (0,-2) {$-2$};
\node[font=\tiny,text=muted,left,inner sep=1.5pt] at (0,-1) {$-1$};
\node[font=\tiny,text=muted,left,inner sep=1.5pt] at (0,1) {$1$};
\node[font=\tiny,text=muted,left,inner sep=1.5pt] at (0,2) {$2$};
\node[font=\tiny,text=muted,left,inner sep=1.5pt] at (0,3) {$3$};
\node[font=\tiny,text=muted,below left,inner sep=1.5pt] at (0,0) {$0$};
\draw[accent,line width=1.1pt] (-0.333,4.000)--(2.333,-4.000);
\end{tikzpicture}\\[1pt]\small d)\end{tabular}\end{center}

\ex{2} Dibuja la recta que pasa por $A(1,\ 2)$ y tiene pendiente 3, y la recta que pasa por $B(0,\ 4)$ y tiene pendiente $-\frac12$.

\ex{2} Una carretera tiene una pendiente del 12\,\%: sube 12\,m por cada 100\,m en horizontal. a) ¿Cuánto sube en 3\,km? b) Una rampa tiene una pendiente del 8\,\% y tiene que subir 0,6\,m. ¿Qué longitud horizontal tiene que tener?

\ex{2} Comprueba con las pendientes si los puntos están alineados:\begin{tasks}(1)
  \task $A(0,\ 1)$, $B(2,\ 7)$, $C(5,\ 16)$
  \task $P(-1,\ 3)$, $Q(1,\ 2)$, $R(5,\ -1)$
\end{tasks}

\ex{2} Sin dibujarlas, ordena las rectas de menos a más inclinada: $r$: pasa por $(0,\ 0)$ y $(2,\ 6)$; \ $s$: pasa por $(1,\ 1)$ y $(5,\ 3)$; \ $t$: pasa por $(0,\ 4)$ y $(1,\ 0)$.

\ex{3} a) ¿Qué valor tiene que tener $k$ para que los puntos $(0,\ 1)$, $(2,\ 5)$ y $(4,\ k)$ estén alineados? b) ¿Qué valor tiene que tener $a$ para que la pendiente de la recta que pasa por $(a,\ 3)$ y $(2,\ 7)$ sea 2?


% =====================================================
\seccio{6}{Funciones lineales: $y=mx+n$}

\ex{1} Escribe la pendiente $m$ y la ordenada en el origen $n$ de cada recta:\begin{tasks}(5)
  \task $y=3x-2$
  \task $y=-x+5$
  \task $y=4-2x$
  \task $y=\frac x2+3$
  \task $y=-6x$
\end{tasks}\recorda{En $y=mx+n$, $m$ es la \textbf{pendiente} y $n$ la \textbf{ordenada en el origen}: la recta corta el eje $Y$ en $(0,\ n)$. Las rectas con la misma pendiente son paralelas.}

\ex{1} Representa en los mismos ejes $y=x+1$ e $y=-x+3$. ¿En qué punto se cortan?

\ex{1} Escribe la ecuación de las rectas $r$, $s$ y $t$:\par\smallskip\begin{center}\begin{tikzpicture}[x=0.42cm,y=0.42cm,baseline=(current bounding box.center),line cap=round,line join=round]
\draw[muted!30,line width=0.3pt] (-4,-4)--(-4,4);
\draw[muted!30,line width=0.3pt] (-3,-4)--(-3,4);
\draw[muted!30,line width=0.3pt] (-2,-4)--(-2,4);
\draw[muted!30,line width=0.3pt] (-1,-4)--(-1,4);
\draw[muted!30,line width=0.3pt] (0,-4)--(0,4);
\draw[muted!30,line width=0.3pt] (1,-4)--(1,4);
\draw[muted!30,line width=0.3pt] (2,-4)--(2,4);
\draw[muted!30,line width=0.3pt] (3,-4)--(3,4);
\draw[muted!30,line width=0.3pt] (4,-4)--(4,4);
\draw[muted!30,line width=0.3pt] (-4,-4)--(4,-4);
\draw[muted!30,line width=0.3pt] (-4,-3)--(4,-3);
\draw[muted!30,line width=0.3pt] (-4,-2)--(4,-2);
\draw[muted!30,line width=0.3pt] (-4,-1)--(4,-1);
\draw[muted!30,line width=0.3pt] (-4,0)--(4,0);
\draw[muted!30,line width=0.3pt] (-4,1)--(4,1);
\draw[muted!30,line width=0.3pt] (-4,2)--(4,2);
\draw[muted!30,line width=0.3pt] (-4,3)--(4,3);
\draw[muted!30,line width=0.3pt] (-4,4)--(4,4);
\draw[-latex,muted,line width=0.6pt] (-4,0)--(4.3,0) node[right,font=\scriptsize,text=navy]{x};
\draw[-latex,muted,line width=0.6pt] (0,-4)--(0,4.3) node[above,font=\scriptsize,text=navy]{y};
\node[font=\tiny,text=muted,below,inner sep=1.5pt] at (-3,0) {$-3$};
\node[font=\tiny,text=muted,below,inner sep=1.5pt] at (-2,0) {$-2$};
\node[font=\tiny,text=muted,below,inner sep=1.5pt] at (-1,0) {$-1$};
\node[font=\tiny,text=muted,below,inner sep=1.5pt] at (1,0) {$1$};
\node[font=\tiny,text=muted,below,inner sep=1.5pt] at (2,0) {$2$};
\node[font=\tiny,text=muted,below,inner sep=1.5pt] at (3,0) {$3$};
\node[font=\tiny,text=muted,left,inner sep=1.5pt] at (0,-3) {$-3$};
\node[font=\tiny,text=muted,left,inner sep=1.5pt] at (0,-2) {$-2$};
\node[font=\tiny,text=muted,left,inner sep=1.5pt] at (0,-1) {$-1$};
\node[font=\tiny,text=muted,left,inner sep=1.5pt] at (0,1) {$1$};
\node[font=\tiny,text=muted,left,inner sep=1.5pt] at (0,2) {$2$};
\node[font=\tiny,text=muted,left,inner sep=1.5pt] at (0,3) {$3$};
\node[font=\tiny,text=muted,below left,inner sep=1.5pt] at (0,0) {$0$};
\draw[accent,line width=1.1pt] (-1.500,-4.000)--(2.500,4.000);
\draw[teal,line width=1.1pt] (-3.000,4.000)--(4.000,-0.667);
\draw[navy,line width=1.1pt] (-4.000,-3.000)--(4.000,-3.000);
\node[font=\scriptsize,right] at (2.2,3.6) {$r$};
\node[font=\scriptsize,below] at (-3.6,3.6) {$s$};
\node[font=\scriptsize,above] at (3.6,-3) {$t$};
\end{tikzpicture}\end{center}

\ex{2} Escribe la ecuación de la recta que:\begin{tasks}(1)
  \task tiene pendiente 3 y pasa por el punto $(1,\ -2)$;
  \task tiene pendiente $-\frac12$ y pasa por el punto $(4,\ 1)$;
  \task es paralela a $y=-2x+1$ y pasa por $(3,\ 0)$.
\end{tasks}

\ex{2} Halla la ecuación de la recta que pasa por cada par de puntos:\begin{tasks}(2)
  \task $(-1,\ 4)$ y $(2,\ -2)$
  \task $(0,\ -3)$ y $(4,\ 3)$
\end{tasks}

\ex{2} Comprueba si los puntos $A(4,\ 7)$, $B(-1,\ -8)$ y $C(0{,}5;\ -3{,}5)$ son de la recta $y=3x-5$.

\ex{3} La recta $y=2x-6$ corta los ejes en dos puntos. Hállalos y calcula el área del triángulo que forma la recta con los ejes.

\ex{3} El agua se congela a 0\,°C, que son 32\,°F (grados Fahrenheit), y hierve a 100\,°C, que son 212\,°F. La relación es lineal. a) Halla la fórmula $F=mC+n$. b) ¿Cuántos °F son 25\,°C? c) ¿Cuántos °C son 50\,°F?


% =====================================================
\seccio{7}{Funciones constantes: $y=k$}

\ex{1} Representa en los mismos ejes las funciones $y=2$, $y=-3$ e $y=0$. ¿Qué pendiente tienen?\recorda{Una \textbf{función constante} $y=k$ vale siempre lo mismo: la gráfica es una recta \textbf{horizontal} y la pendiente es 0. Una recta vertical $x=a$ \textbf{no} es una función.}

\ex{1} Escribe la ecuación de la recta horizontal que pasa por $(4,\ -1)$ y de la recta vertical que pasa por $(-2,\ 5)$. ¿Cuál es una función?

\ex{1} Clasifica las funciones en de proporcionalidad, lineales (con $n\neq0$), constantes o ninguna de estas:\begin{tasks}(4)
  \task $y=-4x$
  \task $y=7$
  \task $y=2x+1$
  \task $y=x^2+1$
  \task $y=\frac{x}{5}$
  \task $y=3-x$
  \task $y=\frac{6}{x}$
  \task $y=-2$
\end{tasks}

\ex{2} Enviar un paquete cuesta 4\,€ hasta 2\,kg, 7\,€ de 2 a 5\,kg y 10\,€ de 5 a 10\,kg. a) Dibuja la gráfica. b) ¿Cuánto cuesta enviar un paquete de 3,5\,kg? c) ¿Qué sale más barato: enviar dos paquetes de 3\,kg o uno de 6\,kg?

\ex{2} Una compañía telefónica ofrece una tarifa plana de 30\,€ al mes y otra cobra 0,08\,€ por minuto. a) Escribe las dos funciones. b) ¿Cuánto pagamos con cada una si hablamos 200 minutos? c) ¿A partir de cuántos minutos conviene la tarifa plana?

\ex{3} La tabla corresponde a una función lineal. Halla su ecuación y explica cómo es la gráfica. ¿Qué valor tendrá $y$ para $x=100$?\par\smallskip\begin{center}$\begin{array}{|c|ccccc|}\hline x&-2&-1&0&1&2\\\hline y&5&5&5&5&5\\\hline\end{array}$\end{center}


% =====================================================
\seccio{8}{Interpretación de gráficas y problemas}

\ex{1} La gráfica muestra la distancia a casa de un excursionista. a) ¿A qué distancia estaba a la 1\,h y media? b) ¿Cuánto tiempo descansó? c) ¿Qué distancia máxima se alejó? d) ¿Cuántos km recorrió en total? e) ¿Va más deprisa a la ida o a la vuelta?\par\smallskip\begin{center}\begin{tikzpicture}[x=0.9cm,y=0.3cm,baseline=(current bounding box.center),line cap=round,line join=round]
\draw[muted!30,line width=0.3pt] (0,0)--(0,11.5);
\draw[muted!30,line width=0.3pt] (1,0)--(1,11.5);
\draw[muted!30,line width=0.3pt] (2,0)--(2,11.5);
\draw[muted!30,line width=0.3pt] (3,0)--(3,11.5);
\draw[muted!30,line width=0.3pt] (4,0)--(4,11.5);
\draw[muted!30,line width=0.3pt] (5,0)--(5,11.5);
\draw[muted!30,line width=0.3pt] (0,0)--(5.6,0);
\draw[muted!30,line width=0.3pt] (0,1)--(5.6,1);
\draw[muted!30,line width=0.3pt] (0,2)--(5.6,2);
\draw[muted!30,line width=0.3pt] (0,3)--(5.6,3);
\draw[muted!30,line width=0.3pt] (0,4)--(5.6,4);
\draw[muted!30,line width=0.3pt] (0,5)--(5.6,5);
\draw[muted!30,line width=0.3pt] (0,6)--(5.6,6);
\draw[muted!30,line width=0.3pt] (0,7)--(5.6,7);
\draw[muted!30,line width=0.3pt] (0,8)--(5.6,8);
\draw[muted!30,line width=0.3pt] (0,9)--(5.6,9);
\draw[muted!30,line width=0.3pt] (0,10)--(5.6,10);
\draw[muted!30,line width=0.3pt] (0,11)--(5.6,11);
\draw[-latex,muted,line width=0.6pt] (0,0)--(5.9,0) node[right,font=\scriptsize,text=navy]{t (h)};
\draw[-latex,muted,line width=0.6pt] (0,0)--(0,11.8) node[above,font=\scriptsize,text=navy]{d (km)};
\node[font=\tiny,text=muted,below,inner sep=1.5pt] at (1,0) {$1$};
\node[font=\tiny,text=muted,below,inner sep=1.5pt] at (2,0) {$2$};
\node[font=\tiny,text=muted,below,inner sep=1.5pt] at (3,0) {$3$};
\node[font=\tiny,text=muted,below,inner sep=1.5pt] at (4,0) {$4$};
\node[font=\tiny,text=muted,below,inner sep=1.5pt] at (5,0) {$5$};
\node[font=\tiny,text=muted,left,inner sep=1.5pt] at (0,2) {$2$};
\node[font=\tiny,text=muted,left,inner sep=1.5pt] at (0,4) {$4$};
\node[font=\tiny,text=muted,left,inner sep=1.5pt] at (0,6) {$6$};
\node[font=\tiny,text=muted,left,inner sep=1.5pt] at (0,8) {$8$};
\node[font=\tiny,text=muted,left,inner sep=1.5pt] at (0,10) {$10$};
\node[font=\tiny,text=muted,below left,inner sep=1.5pt] at (0,0) {$0$};
\draw[accent,line width=1.1pt] (0.000,0.000)--(1.000,6.000)--(2.000,6.000)--(3.000,10.000)--(5.000,0.000);
\fill[navy] (0,0) circle (2.2pt);
\fill[navy] (1,6) circle (2.2pt);
\fill[navy] (2,6) circle (2.2pt);
\fill[navy] (3,10) circle (2.2pt);
\fill[navy] (5,0) circle (2.2pt);
\end{tikzpicture}\end{center}\recorda{Leemos la gráfica de izquierda a derecha: donde sube se aleja, donde es horizontal está parado y donde baja vuelve. Cuanto más inclinada, más deprisa.}

\ex{1} La gráfica da la temperatura de un día. Di dónde es creciente y dónde es decreciente, y cuáles son la temperatura máxima y la mínima y cuándo se alcanzaron.\par\smallskip\begin{center}\begin{tikzpicture}[x=0.3cm,y=0.2cm,baseline=(current bounding box.center),line cap=round,line join=round]
\draw[muted!30,line width=0.3pt] (0,0)--(0,20);
\draw[muted!30,line width=0.3pt] (2,0)--(2,20);
\draw[muted!30,line width=0.3pt] (4,0)--(4,20);
\draw[muted!30,line width=0.3pt] (6,0)--(6,20);
\draw[muted!30,line width=0.3pt] (8,0)--(8,20);
\draw[muted!30,line width=0.3pt] (10,0)--(10,20);
\draw[muted!30,line width=0.3pt] (12,0)--(12,20);
\draw[muted!30,line width=0.3pt] (14,0)--(14,20);
\draw[muted!30,line width=0.3pt] (16,0)--(16,20);
\draw[muted!30,line width=0.3pt] (18,0)--(18,20);
\draw[muted!30,line width=0.3pt] (20,0)--(20,20);
\draw[muted!30,line width=0.3pt] (22,0)--(22,20);
\draw[muted!30,line width=0.3pt] (24,0)--(24,20);
\draw[muted!30,line width=0.3pt] (0,0)--(25,0);
\draw[muted!30,line width=0.3pt] (0,2)--(25,2);
\draw[muted!30,line width=0.3pt] (0,4)--(25,4);
\draw[muted!30,line width=0.3pt] (0,6)--(25,6);
\draw[muted!30,line width=0.3pt] (0,8)--(25,8);
\draw[muted!30,line width=0.3pt] (0,10)--(25,10);
\draw[muted!30,line width=0.3pt] (0,12)--(25,12);
\draw[muted!30,line width=0.3pt] (0,14)--(25,14);
\draw[muted!30,line width=0.3pt] (0,16)--(25,16);
\draw[muted!30,line width=0.3pt] (0,18)--(25,18);
\draw[muted!30,line width=0.3pt] (0,20)--(25,20);
\draw[-latex,muted,line width=0.6pt] (0,0)--(25.6,0) node[right,font=\scriptsize,text=navy]{h};
\draw[-latex,muted,line width=0.6pt] (0,0)--(0,20.6) node[above,font=\scriptsize,text=navy]{°C};
\node[font=\tiny,text=muted,below,inner sep=1.5pt] at (4,0) {$4$};
\node[font=\tiny,text=muted,below,inner sep=1.5pt] at (8,0) {$8$};
\node[font=\tiny,text=muted,below,inner sep=1.5pt] at (12,0) {$12$};
\node[font=\tiny,text=muted,below,inner sep=1.5pt] at (16,0) {$16$};
\node[font=\tiny,text=muted,below,inner sep=1.5pt] at (20,0) {$20$};
\node[font=\tiny,text=muted,below,inner sep=1.5pt] at (24,0) {$24$};
\node[font=\tiny,text=muted,left,inner sep=1.5pt] at (0,4) {$4$};
\node[font=\tiny,text=muted,left,inner sep=1.5pt] at (0,8) {$8$};
\node[font=\tiny,text=muted,left,inner sep=1.5pt] at (0,12) {$12$};
\node[font=\tiny,text=muted,left,inner sep=1.5pt] at (0,16) {$16$};
\node[font=\tiny,text=muted,below left,inner sep=1.5pt] at (0,0) {$0$};
\draw[accent,line width=1.1pt] (0.000,10.000)--(4.000,6.000)--(8.000,8.000)--(12.000,16.000)--(16.000,18.000)--(20.000,12.000)--(24.000,9.000);
\fill[navy] (0,10) circle (2.2pt);
\fill[navy] (4,6) circle (2.2pt);
\fill[navy] (8,8) circle (2.2pt);
\fill[navy] (12,16) circle (2.2pt);
\fill[navy] (16,18) circle (2.2pt);
\fill[navy] (20,12) circle (2.2pt);
\fill[navy] (24,9) circle (2.2pt);
\end{tikzpicture}\end{center}

\ex{2} Relaciona cada historia con su gráfica (tiempo, distancia a casa):\begin{tasks}(1)
  \task Voy en bici a casa de una amiga, me quedo un rato y vuelvo.
  \task Salgo corriendo, paro a descansar y después sigo andando despacio.
  \task Camino siempre a la misma velocidad, sin parar.
\end{tasks}\par\smallskip\begin{center}\begin{tabular}{c}\begin{tikzpicture}[x=0.28cm,y=0.28cm,baseline=(current bounding box.center),line cap=round,line join=round]
\draw[muted!30,line width=0.3pt] (0,0)--(0,6.5);
\draw[muted!30,line width=0.3pt] (2,0)--(2,6.5);
\draw[muted!30,line width=0.3pt] (4,0)--(4,6.5);
\draw[muted!30,line width=0.3pt] (6,0)--(6,6.5);
\draw[muted!30,line width=0.3pt] (0,0)--(6.5,0);
\draw[muted!30,line width=0.3pt] (0,2)--(6.5,2);
\draw[muted!30,line width=0.3pt] (0,4)--(6.5,4);
\draw[muted!30,line width=0.3pt] (0,6)--(6.5,6);
\draw[-latex,muted,line width=0.6pt] (0,0)--(7.1,0) node[right,font=\scriptsize,text=navy]{t};
\draw[-latex,muted,line width=0.6pt] (0,0)--(0,7.1) node[above,font=\scriptsize,text=navy]{d};
\node[font=\tiny,text=muted,below,inner sep=1.5pt] at (2,0) {$2$};
\node[font=\tiny,text=muted,below,inner sep=1.5pt] at (4,0) {$4$};
\node[font=\tiny,text=muted,below,inner sep=1.5pt] at (6,0) {$6$};
\node[font=\tiny,text=muted,left,inner sep=1.5pt] at (0,2) {$2$};
\node[font=\tiny,text=muted,left,inner sep=1.5pt] at (0,4) {$4$};
\node[font=\tiny,text=muted,left,inner sep=1.5pt] at (0,6) {$6$};
\node[font=\tiny,text=muted,below left,inner sep=1.5pt] at (0,0) {$0$};
\draw[accent,line width=1.1pt] (0.000,0.000)--(2.000,5.000)--(4.000,5.000)--(6.000,0.000);
\end{tikzpicture}\\[1pt]\small A\end{tabular}\hspace{0.7cm}\begin{tabular}{c}\begin{tikzpicture}[x=0.28cm,y=0.28cm,baseline=(current bounding box.center),line cap=round,line join=round]
\draw[muted!30,line width=0.3pt] (0,0)--(0,6.5);
\draw[muted!30,line width=0.3pt] (2,0)--(2,6.5);
\draw[muted!30,line width=0.3pt] (4,0)--(4,6.5);
\draw[muted!30,line width=0.3pt] (6,0)--(6,6.5);
\draw[muted!30,line width=0.3pt] (0,0)--(6.5,0);
\draw[muted!30,line width=0.3pt] (0,2)--(6.5,2);
\draw[muted!30,line width=0.3pt] (0,4)--(6.5,4);
\draw[muted!30,line width=0.3pt] (0,6)--(6.5,6);
\draw[-latex,muted,line width=0.6pt] (0,0)--(7.1,0) node[right,font=\scriptsize,text=navy]{t};
\draw[-latex,muted,line width=0.6pt] (0,0)--(0,7.1) node[above,font=\scriptsize,text=navy]{d};
\node[font=\tiny,text=muted,below,inner sep=1.5pt] at (2,0) {$2$};
\node[font=\tiny,text=muted,below,inner sep=1.5pt] at (4,0) {$4$};
\node[font=\tiny,text=muted,below,inner sep=1.5pt] at (6,0) {$6$};
\node[font=\tiny,text=muted,left,inner sep=1.5pt] at (0,2) {$2$};
\node[font=\tiny,text=muted,left,inner sep=1.5pt] at (0,4) {$4$};
\node[font=\tiny,text=muted,left,inner sep=1.5pt] at (0,6) {$6$};
\node[font=\tiny,text=muted,below left,inner sep=1.5pt] at (0,0) {$0$};
\draw[accent,line width=1.1pt] (0.000,0.000)--(3.000,3.000)--(6.000,6.000);
\end{tikzpicture}\\[1pt]\small B\end{tabular}\hspace{0.7cm}\begin{tabular}{c}\begin{tikzpicture}[x=0.28cm,y=0.28cm,baseline=(current bounding box.center),line cap=round,line join=round]
\draw[muted!30,line width=0.3pt] (0,0)--(0,6.5);
\draw[muted!30,line width=0.3pt] (2,0)--(2,6.5);
\draw[muted!30,line width=0.3pt] (4,0)--(4,6.5);
\draw[muted!30,line width=0.3pt] (6,0)--(6,6.5);
\draw[muted!30,line width=0.3pt] (0,0)--(6.5,0);
\draw[muted!30,line width=0.3pt] (0,2)--(6.5,2);
\draw[muted!30,line width=0.3pt] (0,4)--(6.5,4);
\draw[muted!30,line width=0.3pt] (0,6)--(6.5,6);
\draw[-latex,muted,line width=0.6pt] (0,0)--(7.1,0) node[right,font=\scriptsize,text=navy]{t};
\draw[-latex,muted,line width=0.6pt] (0,0)--(0,7.1) node[above,font=\scriptsize,text=navy]{d};
\node[font=\tiny,text=muted,below,inner sep=1.5pt] at (2,0) {$2$};
\node[font=\tiny,text=muted,below,inner sep=1.5pt] at (4,0) {$4$};
\node[font=\tiny,text=muted,below,inner sep=1.5pt] at (6,0) {$6$};
\node[font=\tiny,text=muted,left,inner sep=1.5pt] at (0,2) {$2$};
\node[font=\tiny,text=muted,left,inner sep=1.5pt] at (0,4) {$4$};
\node[font=\tiny,text=muted,left,inner sep=1.5pt] at (0,6) {$6$};
\node[font=\tiny,text=muted,below left,inner sep=1.5pt] at (0,0) {$0$};
\draw[accent,line width=1.1pt] (0.000,0.000)--(1.000,5.000)--(3.000,5.000)--(6.000,6.000);
\end{tikzpicture}\\[1pt]\small C\end{tabular}\end{center}

\ex{2} Un taxi cobra 3,50\,€ por la bajada de bandera y 1,10\,€ por km. a) Escribe la función del precio. b) ¿Cuánto cuesta un viaje de 12\,km? c) ¿Cuántos km hemos hecho si hemos pagado 25,50\,€?\recorda{Muchas situaciones tienen una parte fija $n$ y una parte que aumenta (o disminuye) siempre igual por cada unidad, $m$: son funciones lineales $y=mx+n$.}

\ex{2} Un depósito tiene 800 litros y se vacían 25 cada minuto. a) Escribe la función de los litros que quedan. b) ¿Cuántos litros quedan a los 12 minutos? c) ¿Cuándo se vaciará?

\ex{2} Clara sale de casa a las 10\,h y camina 3\,km en 45 minutos hasta la piscina. Está allí 1\,h 30\,min y vuelve en autobús en 15 minutos. Dibuja la gráfica (hora, distancia a casa).

\ex{3} El gimnasio A cobra 30\,€ de matrícula y 15\,€ al mes, y el gimnasio B, 45\,€ de matrícula y 10\,€ al mes. a) Escribe las dos funciones. b) ¿Cuánto cuesta cada uno si vamos 6 meses? c) ¿Para cuántos meses cuestan lo mismo?

\ex{3} Un ciclista sale de un pueblo a 18\,km/h. Una hora después sale otro por el mismo camino a 24\,km/h. a) Escribe la distancia recorrida por cada uno en función del tiempo $t$ que hace que ha salido el primero. b) ¿Cuándo y dónde lo alcanza?

\ex{3} Encendemos a la vez dos velas: una mide 30\,cm y se quema 2\,cm cada hora; la otra mide 24\,cm y se quema 1,5\,cm cada hora. a) Escribe las dos funciones. b) ¿Cuándo tendrán la misma altura? ¿Cuál? c) ¿Cuál se acaba antes?

\vfill
\begin{center}\color{muted}\small\faCheckCircle[regular]\ \ ¡Ánimo! Comprueba tus resultados con las herramientas interactivas de los apuntes de la unidad.\end{center}

\end{document}
